\section*{Abstract}
Cloud computing pipelines, particularly Continuous Integration and Continuous Deployment (CI/CD) environments, are essential for modern software development due to their flexibility, scalability, and automation. However, these dynamic cloud infrastructures face continuously evolving cyber threats, including polymorphic and metamorphic malware, which traditional static detection systems struggle to handle. Existing systems rely on fixed architectures, require new threat samples for retraining, and lack self-learning capabilities in sync with the pipeline, often resulting in high false positives, operational delays, and potential data breaches or financial loss. To address these limitations, this research proposes an adaptive malware detection system for cloud CI/CD pipelines. The system uses deep neural networks to perform both static and behavioral analysis of pipeline artifacts such as code, build logs, and containers. A key feature is its self-evolving capability, where the neural network can restructure and adjust its parameters dynamically in response to new behaviors, enabling the system to adapt to previously unseen threats. Unlike post-deployment scanners, this model is embedded directly in the pipeline, allowing real-time monitoring and detection. Additionally, the system includes a threat mutation module that generates modified versions of known malware and injects them into the detection environment. This allows the model to learn from controlled adversarial samples while distinguishing them from external threats, thereby improving detection accuracy and reducing false positives. By continuously learning from both normal operational patterns and self-generated threats, the system maintains resilience against evolving malware and strengthens the security of cloud pipelines. The expected contribution of this research is a self-adaptive neural framework capable of evolving with changing threat behavior and providing proactive, context-aware protection for cloud CI/CD pipelines, addressing the shortcomings of current static detection systems, and enhancing the reliability, efficiency, and security of modern cloud software delivery.


\vspace{1cm}
\textbf{Keywords:} Adaptive Evolving Neural Security Architectures, Cloud-Native CI/CD Security and Real-Time Defense, Polymorphic and Metamorphic Malware Detection, Behavioral and Code-Level Anomaly Detection, Static and Dynamic Threat Analysis (SAST/DAST), Adversarial Training, Red Team Simulation
\pagebreak
